\documentclass[aspectratio=169,11pt]{beamer}
\usepackage[T1]{fontenc}
\usepackage{lmodern}
\usepackage{listings}
\usepackage{booktabs}
\definecolor{navy}{RGB}{25,48,76}
\definecolor{teal}{RGB}{0,112,120}
\definecolor{codebg}{RGB}{246,247,249}
\setbeamercolor{structure}{fg=teal}
\setbeamercolor{frametitle}{fg=navy}
\setbeamercolor{title}{fg=navy}
\setbeamercolor{normal text}{fg=navy,bg=white}
\setbeamertemplate{navigation symbols}{}
\setbeamertemplate{footline}{\hfill\usebeamercolor[fg]{structure}\insertframenumber\,/\,\inserttotalframenumber\hspace{5mm}\vspace{3mm}}
\setbeamertemplate{itemize items}[circle]
\setbeamersize{text margin left=9mm,text margin right=9mm}
\lstset{language=Python,basicstyle=\ttfamily\small,keywordstyle=\color{teal}\bfseries,commentstyle=\color{gray},stringstyle=\color{teal},backgroundcolor=\color{codebg},frame=none,showstringspaces=false,columns=fullflexible,keepspaces=true,breaklines=true,xleftmargin=2mm,xrightmargin=2mm,aboveskip=8pt,belowskip=8pt}
\newcommand{\code}[1]{\texttt{#1}}
\usepackage{amsmath}
\title{Testing Python with unittest}
\subtitle{Test cases, fixtures, and controlled dependencies}
\author{CS413}
\date{October 8, 2026}
\begin{document}
\begin{frame}\titlepage\end{frame}
\begin{frame}[fragile]{Python's built-in testing framework}
\begin{itemize}
\item \code{unittest} comes with Python. No separate installation is needed.
\item Tests are methods on subclasses of \code{unittest.TestCase}.
\item A runner executes the tests and reports their outcomes.
\end{itemize}
\begin{lstlisting}
python3 -m unittest -v
\end{lstlisting}
\end{frame}

\begin{frame}[fragile]{An executable contract}
\begin{itemize}
\item For a stack, the last item pushed should be the first item popped.
\item Popping an empty stack should raise \code{IndexError}.
\item Tests turn these expectations into repeatable checks.
\item Passing examples provide evidence. They do not prove correctness for every input.
\end{itemize}
\end{frame}

\begin{frame}[fragile]{The stack under test}
\begin{lstlisting}
# stack.py
class Stack:
    def __init__(self):
        self._items = []

    def push(self, value):
        self._items.append(value)

    def pop(self):
        if not self._items:
            raise IndexError("empty stack")
        return self._items.pop()
\end{lstlisting}
\end{frame}

\begin{frame}[fragile]{A first TestCase}
\begin{lstlisting}
# test_stack.py
import unittest
from stack import Stack

class TestStack(unittest.TestCase):
    def test_pop_returns_last_pushed(self):
        stack = Stack()
        stack.push("x")
        stack.push("y")
        self.assertEqual(stack.pop(), "y")

if __name__ == "__main__":
    unittest.main()
\end{lstlisting}
The guard enables direct execution without running the suite on import.
\end{frame}

\begin{frame}[fragile]{Running the examples}
\begin{lstlisting}
cd examples/unittest_demo
python3 -m unittest -v
python3 -m unittest test_stack
python3 -m unittest test_stack.TestStack
python3 -m unittest test_stack.TestStack.test_empty_pop
python3 test_stack.py -v
\end{lstlisting}
The last command uses the \code{unittest.main()} guard.

All example commands run from this demo directory.
\end{frame}

\begin{frame}[fragile]{Discovery and selection}
\begin{lstlisting}
python3 -m unittest discover -s . -p 'test*.py' -v
python3 -m unittest -k pop
python3 -m unittest -f
python3 -m unittest -b
\end{lstlisting}
\begin{itemize}
\item Default discovery uses \code{test*.py} files and methods starting with \code{test}.
\item \code{-k} matches a substring or wildcard pattern, not a pytest Boolean expression.
\item \code{-f} stops at the first failure or error. \code{-b} buffers output.
\item For nested test directories, use importable packages with \code{\_\_init\_\_.py}.
\end{itemize}
\end{frame}

\begin{frame}[fragile]{Failures and errors}
Suppose \code{pop()} incorrectly calls \code{self.\_items.pop(0)}.
\begin{lstlisting}
FAIL: test_pop_returns_last_pushed (...)
...
AssertionError: 'x' != 'y'
- x
+ y
\end{lstlisting}
\begin{itemize}
\item A failed assertion reports a failure.
\item An unexpected exception, such as \code{NameError}, reports an error.
\item Read the failing expectation and traceback before changing the code.
\end{itemize}
\end{frame}

\begin{frame}[fragile]{Assertion methods express the expectation}
\renewcommand{\arraystretch}{1.4}
\begin{tabular}{ll}
\toprule
Method & Meaning \\
\midrule
\code{assertEqual(a, b)} & Equal values \\
\code{assertIs(a, b)} & The same object \\
\code{assertIsNone(x)} & The value is \code{None} \\
\code{assertIn(x, xs)} & Membership \\
\code{assertTrue(p)} & Truthiness \\
\code{assertCountEqual(a, b)} & Same elements and multiplicities \\
\bottomrule
\end{tabular}
\medskip

For equality, \code{assertEqual(a, b)} usually gives a clearer diagnostic
than \code{assertTrue(a == b)}.
\end{frame}

\begin{frame}[fragile]{Values, identity, and order}
\begin{lstlisting}
def test_list_behavior(self):
    a = [1, 2]
    b = [1, 2]
    self.assertEqual(a, b)
    self.assertIsNot(a, b)
    self.assertCountEqual([1, 2, 2], [2, 1, 2])
\end{lstlisting}
\begin{itemize}
\item Choose the assertion from the contract.
\item A stack test needs order-sensitive equality.
\item Use \code{assertCountEqual} only when ordering is irrelevant.
\end{itemize}
\end{frame}

\begin{frame}[fragile]{Expected exceptions}
\begin{lstlisting}
def test_empty_pop(self):
    stack = Stack()
    with self.assertRaisesRegex(IndexError, "empty stack"):
        stack.pop()
\end{lstlisting}
\begin{itemize}
\item \code{assertRaises} checks the exception type.
\item \code{assertRaisesRegex} also searches the message with a regular expression.
\item No exception or a different exception makes the test unsuccessful.
\item Keep unrelated setup outside the context manager.
\end{itemize}
\end{frame}

\begin{frame}[fragile]{Fresh state with setUp}
\begin{lstlisting}
class TestStack(unittest.TestCase):
    def setUp(self):
        self.stack = Stack()

    def test_round_trip(self):
        self.stack.push(7)
        self.assertEqual(self.stack.pop(), 7)

    def test_empty_pop(self):
        with self.assertRaises(IndexError):
            self.stack.pop()
\end{lstlisting}
Each test gets a new \code{TestCase} instance and its own \code{setUp()} call.
\end{frame}

\begin{frame}[fragile]{Fixture lifecycle}
\begin{itemize}
\item For a normal test: \code{setUp()}, test method, \code{tearDown()}, cleanups.
\item \code{tearDown()} runs after a failing test if \code{setUp()} succeeded.
\item If \code{setUp()} raises, the test and \code{tearDown()} do not run.
\item Registered \code{addCleanup} callbacks still run if setup fails.
\item Cleanup callbacks run in reverse registration order.
\end{itemize}
\end{frame}

\begin{frame}[fragile]{Temporary resources and cleanup}
\begin{lstlisting}
import tempfile
from pathlib import Path

class TestFiles(unittest.TestCase):
    def setUp(self):
        temp = tempfile.TemporaryDirectory()
        self.addCleanup(temp.cleanup)
        self.directory = Path(temp.name)

    def test_text_round_trip(self):
        path = self.directory / "out.txt"
        path.write_text("hello", encoding="utf-8")
        self.assertEqual(path.read_text(encoding="utf-8"),
                         "hello")
\end{lstlisting}
Register cleanup immediately after acquiring the resource.
\end{frame}

\begin{frame}[fragile]{Shared fixtures and test independence}
\begin{itemize}
\item \code{setUpClass} and \code{tearDownClass} manage class-level resources.
\item Use \code{@classmethod} for these methods.
\item A shared mutable resource can leak state between tests.
\item Prefer per-test setup until the cost justifies sharing.
\item Each test should pass when run alone. Tests must not depend on their execution order.
\end{itemize}
\end{frame}

\begin{frame}[fragile]{Multiple cases with subTest}
\begin{lstlisting}
def test_lifo_cases(self):
    for values in ([], [7], [2, 5, 2]):
        with self.subTest(values=values):
            stack = Stack()
            for value in values:
                stack.push(value)
            actual = [stack.pop() for _ in values]
            self.assertEqual(actual, list(reversed(values)))
\end{lstlisting}
\begin{itemize}
\item A failing subtest reports its parameters and normally allows the loop to continue.
\item These are subcases of one test method, not separately discovered methods.
\item Construct fresh state inside each subtest.
\end{itemize}
\end{frame}

\begin{frame}[fragile]{Choosing useful cases}
\begin{itemize}
\item An empty stack exercises the error contract.
\item Distinct values expose FIFO behavior accidentally implemented as LIFO.
\item Repeated values detect implementations that discard duplicates.
\item Interleaved pushes and pops exercise changes in state.
\item A regression test should fail with the specific bug and pass after the fix.
\end{itemize}
\end{frame}

\begin{frame}[fragile]{Floating-point expectations}
\begin{lstlisting}
def test_decimal_sum(self):
    self.assertAlmostEqual(0.1 + 0.2, 0.3, delta=1e-12)
\end{lstlisting}
\begin{itemize}
\item \code{delta} sets an absolute tolerance.
\item Choose the tolerance from the numerical requirements.
\item The default \code{places=7} checks rounded differences, not significant digits.
\item For a relative tolerance, use \code{math.isclose} with explicit tolerances.
\end{itemize}
\end{frame}

\begin{frame}[fragile]{Capturing printed output}
\begin{lstlisting}
from contextlib import redirect_stdout
from io import StringIO
from stack import announce

def test_announce(self):
    output = StringIO()
    with redirect_stdout(output):
        announce("Ada")
    self.assertEqual(output.getvalue(), "Hello, Ada!\n")
\end{lstlisting}
Here \code{announce(name)} calls \code{print(f"Hello, \{name\}!")}.

The newline is part of the expected output.
\end{frame}

\begin{frame}[fragile]{Controlling environment variables}
\begin{lstlisting}
# stack.py
import os

def greeting():
    return "Hello, " + os.getenv("USER_NAME", "guest")

# Inside a TestCase
from unittest.mock import patch

def test_greeting(self):
    with patch.dict("os.environ", {"USER_NAME": "Ada"}):
        self.assertEqual(greeting(), "Hello, Ada")
\end{lstlisting}
Place the imports at module scope. The context restores the environment.
\end{frame}

\begin{frame}[fragile]{Mocking an external dependency}
\begin{lstlisting}
# service.py
from urllib.request import urlopen

def status_code(url):
    with urlopen(url, timeout=2) as response:
        return response.status

# Inside a TestCase; patch is imported at module scope
@patch("service.urlopen", autospec=True)
def test_status_code(self, mock_open):
    mock_open.return_value.__enter__.return_value.status = 200
    self.assertEqual(status_code("https://example.org"), 200)
    mock_open.assert_called_once_with("https://example.org",
                                      timeout=2)
\end{lstlisting}
\end{frame}

\begin{frame}[fragile]{Where to patch, and what a mock proves}
\begin{itemize}
\item Patch \code{service.urlopen}: that is where \code{status\_code} looks up the name.
\item Patching the original library after \code{service} imports the name would miss this reference.
\item \code{autospec=True} constrains the mock to the original interface where inspectable.
\item The test checks our interaction with a substitute. It does not establish that the network service works.
\item Separate integration tests can exercise the real dependency.
\end{itemize}
\end{frame}

\begin{frame}[fragile]{Skipped tests and expected failures}
\begin{lstlisting}
@unittest.skip("requires the course server")
def test_server(self):
    ...

@unittest.expectedFailure
def test_known_bug(self):
    self.assertEqual(buggy_function(), expected)
\end{lstlisting}
\begin{itemize}
\item Illustrative placeholders, not part of the runnable demo.
\item A skipped test does not run. An expected failure documents a known bug.
\item An unexpected success makes the unittest run unsuccessful.
\item Use a reason or issue reference so the marker can be revisited.
\end{itemize}
\end{frame}

\begin{frame}[fragile]{unittest and pytest: equivalent intentions}
\renewcommand{\arraystretch}{1.55}
\begin{tabular}{p{.45\textwidth}p{.45\textwidth}}
\toprule
\code{unittest} & \code{pytest} \\
\midrule
\code{TestCase} method & Usually a plain function \\
\code{self.assertEqual(a, b)} & \code{assert a == b} \\
\code{self.assertRaises(E)} & \code{pytest.raises(E)} \\
\code{setUp} / cleanups & Fixtures / yield cleanup \\
\code{subTest} loop & \code{mark.parametrize} cases \\
Standard library & Separately installed package \\
\bottomrule
\end{tabular}
\medskip

Both can express the same behavioral expectations. The setup and
reporting conventions differ.
\end{frame}

\begin{frame}[fragile]{Exercise: free variables}
For tuple-based expressions, the intended rules are:
\[
FV(x)=\{x\},\qquad FV(e_1\ e_2)=FV(e_1)\cup FV(e_2)
\]
\[
FV(\lambda x.e)=FV(e)\setminus\{x\}
\]
\begin{lstlisting}
("var", "x")
("app", ("var", "x"), ("var", "y"))
("lam", "x", ("var", "x"))
\end{lstlisting}
\begin{itemize}
\item Write tests for a bound variable and a variable that remains free.
\item Add an application and a nested lambda, using \code{subTest}.
\item Which case detects a missing subtraction in the lambda rule?
\end{itemize}
\end{frame}

\begin{frame}[fragile]{Exercise discussion}
\begin{lstlisting}
from expr import free_vars

class TestFreeVars(unittest.TestCase):
    def test_cases(self):
        cases = [
            (("lam", "x", ("var", "x")), set()),
            (("lam", "x", ("var", "y")), {"y"}),
            (("app", ("var", "x"), ("var", "y")), {"x", "y"}),
        ]
        for expr, expected in cases:
            with self.subTest(expr=expr):
                self.assertEqual(free_vars(expr), expected)
\end{lstlisting}
The first case detects the missing subtraction. The demo includes a nested case.
\end{frame}

\begin{frame}[fragile]{Testing workflow and references}
\begin{itemize}
\item Write a specific expectation and observe the intended failure.
\item Fix the implementation, rerun the test, then run the suite.
\item Check that the reported test count is plausible. An empty suite proves nothing.
\item Demo: \code{examples/unittest\_demo/}.
\end{itemize}
\medskip
Official Python documentation:
\begin{itemize}
\item \href{https://docs.python.org/3/library/unittest.html}{unittest: cases, assertions, fixtures, and discovery}
\item \href{https://docs.python.org/3/library/unittest.mock.html}{unittest.mock: patching and test doubles}
\end{itemize}
\end{frame}
\end{document}
